% =============================================================================
% cap2_hardware/04_energia_conectividade.tex
% Seção 2.3 -- Energia, armazenamento e conectividade (v2.2 — segunda compactação)
% =============================================================================
\section{Energia, armazenamento e conectividade}
\label{sec:energia-conectividade}

Energia disponível e canal de saída dos dados são, no Campus 2, condições materiais que precedem qualquer escolha algorítmica. A Seção~\ref{sec:caracterizacao-pontos} já evidenciou a assimetria estrutural dos dez pontos: apenas três dispõem de tomada regular, a maioria está exposta a longa insolação direta, e a rede sem fio institucional está nominalmente disponível, mas operacionalmente sujeita a autorizações e janelas de manutenção fora do alcance do projeto. Essas condições orientam o subsistema para o regime \emph{offline-first}, em que a continuidade da aquisição não depende de rede e a sincronização é tratada como camada opcional. Para o pipeline do Capítulo~\ref{cap:pipeline}, isso implica projeto predominantemente em \textbf{lote}, sobre quadros consolidados em servidor após ciclos de coleta.

A discussão energética articula três regimes. Nos pontos com tomada elétrica, câmeras IP comerciais modificadas para operação local e SBC com NPU operam confortavelmente em ciclo contínuo, com folga para iluminação ativa por NIR e janelas de inferência prolongadas, salvo riscos pontuais de queda de tensão amortecidos por fontes estabilizadas. Esse regime define o teto operacional a partir do qual as etapas mais exigentes do pipeline podem migrar para a borda, mas alcança parcela minoritária dos pontos. O regime majoritário é o de \textbf{alimentação fotovoltaica}: SBCs de classe Raspberry com aceleradores neurais dissipam potência não desprezível em operação contínua, exigindo painel, controlador de carga e banco de bateria dimensionados com folga para o ciclo diário e reserva para dias nublados consecutivos --- requisito já reportado como dimensionante em sistemas similares~\cite{velasco2024smartcameratraps,pishdast2025smartcameratraps,sargent2021raspberrypi}. A iluminação ativa por NIR é item energético sensível: emissores potentes resolvem o problema óptico noturno mas inflacionam o consumo médio, deslocando o ponto de operação para fora do envelope viável de painéis e baterias razoáveis. Por esse motivo, reduzir o nível de processamento do SBC para o ESP32-CAM, ou aceitar perda parcial de cobertura noturna, é alternativa legítima em vez de regressão técnica; plataformas embarcadas mais leves oferecem flexibilidade adicional~\cite{tiddeman2023lowpowerhardware}. O terceiro regime, das câmeras alimentadas por \textbf{baterias} (\emph{trail cam} comercial, Seção~\ref{sec:tipologias-arquiteturais}), opera com consumo médio próximo de zero pela combinação de PIR físico e \emph{deep sleep}, ao custo de processamento de borda essencialmente nulo e de cadência de troca ou recarga das baterias incorporada ao protocolo de campo. Consequentemente, o pipeline não pode pressupor processamento de borda uniforme: a divisão borda-servidor varia com o orçamento energético local, e o software deve tolerar essa heterogeneidade.

O armazenamento local, no regime \emph{offline-first}, é canal primário de persistência. Cartões \emph{microSD} de classe industrial, com tolerância térmica ampliada e algoritmos robustos de \emph{wear leveling}, são o suporte usual; cartões de consumo têm histórico documentado de falha precoce sob ciclos intensos de escrita e variação térmica externa~\cite{tiddeman2023lowpowerhardware,sargent2021raspberrypi}. A política de gravação equilibra densidade de informação retida (taxa de quadros, resolução, metadados de inferência) e envelhecimento do meio físico, com rotação programada e marcação de blocos consolidados, definindo a taxa de chegada de dados e o horizonte temporal das análises de re-ID. Operacionalmente, o ciclo logístico é independente de rede: cartões são retirados em campo pelos cuidadores ou pela equipe técnica e transferidos ao servidor (Seção~\ref{sec:decisao-arquitetural}), com cadência semanal ou quinzenal compatível com a rotina da AEX Gatosdoc2. A queda de tensão durante escrita em curso é mitigada por sistemas de arquivos com \emph{journaling} e escrita em blocos atômicos com renomeação final~\cite{sargent2021raspberrypi}.

A conectividade, ao contrário de energia e armazenamento, é variável institucional. A rede sem fio da Universidade está nominalmente disponível, mas seu uso por dispositivos embarcados depende de autorização da Superintendência de Tecnologia da Informação, credenciais compatíveis com autenticação corporativa (tipicamente \emph{802.1X EAP-PEAP}) e configuração consistente em cada nó --- tarefas que se aproximam mais da escala de complexidade burocrática do que da técnica. A rede celular \emph{4G/LTE} opera sob lógica análoga: tecnicamente trivial, institucionalmente custosa. A consequência metodológica é tratar conectividade como camada opcional adicionada sobre um sistema já completo em modo \emph{offline-first}. Quando autorizada, habilita sincronização oportunista, supervisão remota e atualização de modelos em borda; quando indisponível, o sistema continua a operar com cadência de análise atrelada ao ciclo logístico. Nenhuma funcionalidade central do pipeline pode depender de conectividade ininterrupta: validação humana no laço, atualização de modelos e sincronização de resultados são camadas adicionais, não pré-requisitos.
